\section{Introduction}
The KV cache stores the history of a language-model interaction as reusable inference state. A party that observes this state may recover information about the prompt even without access to the client connection. Cache offloading and reuse therefore raise a confidentiality question alongside their performance benefits. Block-oriented management, exemplified by PagedAttention~\cite{paged}, makes small cache blocks a natural unit for both protection and analysis.

\scheme{}~\cite{kvcloak} addresses this question with reversible obfuscation. Its implementation combines a secret column transform, a sparse additive marker, a fresh row permutation for each block, and a secret row transform. The column transform can be folded into attention weights; the other operations protect stored blocks. The design's security argument assigns an important role to the $b!$ possible row permutations. In the published paper, fresh shuffling is the step intended to defeat recovery of a fixed linear transform through chosen plaintexts. This suggests a focused question: does the number of permutations describe the uncertainty that the attacker actually observes?

Our answer is negative for the tested construction. A permutation changes the ordering of a block, but cannot change a block whose rows are all identical. Such blocks are reachable without selecting arbitrary hidden vectors. In the first Value layer, each row is a deterministic function of its input token under the model architectures we evaluate. Repeated characters can tokenize to a long run of the same token. An attacker submits the text through the normal tokenizer and observes its protected cache. Mandatory prefixes are retained; the attacker selects complete blocks inside the repeat run after observing the cache.

The resulting protected blocks vary only with the shuffled position of the additive marker. For a single-row marker, there are $b$ ideal states, rather than $b!$. Their pairwise differences have rank one. Those differences expose the marker direction, and the geometry of the states exposes the columns of the reused row transform. A known probe vector then identifies the implementation's independent two-dimensional column rotations, subject to explicit nondegeneracy conditions. The attacker can apply the recovered parameters to unrelated blocks from the same secret epoch and match the resulting vectors against a dictionary computed from the public, original model.

This is a confidentiality attack under a cache-observation and chosen-plaintext threat model. It is not a method for obtaining cache access from an ordinary text API. Nor does it reconstruct the order of tokens within each shuffled block. Its output is a token multiset for each block. This distinction matters: strong token-content leakage is possible even when the permutation continues to hide ordering. We evaluate these properties separately rather than equating multiset overlap with exact prompt reconstruction.

The evidence spans several complementary settings. Our main suite uses ten models and twelve configurations, including tokenizer prefixes, different head dimensions, and block sizes 8, 16, and 32. It records raw correct/total counts rather than reconstructing them from rounded rates. A separate test attacks cache files produced by the official command-line entry point for 16 LMSYS conversations. An isolated fused-mode experiment constructs the attack dictionary from the original public model, not the defender's fused weights. A further three-model study varies complete secret configurations over five epochs per model. These experiments support the mechanism while leaving deployment prevalence and unseen-conversation generalization unmeasured.

We also examine what the recovered contents mean for privacy and mitigation. A closed-set synthetic study ranks 100 same-type value candidates in each of 30 cases. The true candidate always attains the maximum score, but ties remain: only 21 cases have a unique maximum. This is a bounded observation, with a fixed candidate-selection protocol and incomplete retained traces, not a general identification-rate claim. Finally, a request-level refresh of the row transform and marker prevents the existing recovered demixer from transferring in one measured setting. This simple intervention also narrows the research conclusion: the attack concerns reused structural secrets, and a new obfuscation system is not needed to state that finding.

Our contributions are:
\begin{itemize}
\item We characterize the observable orbit induced by chosen-input symmetry in shuffled cache obfuscation, deriving state counts, identification conditions, and a collection bound. The analysis distinguishes the single-marker result from more general marker patterns.
\item We construct a repeated-token attack on the official \scheme{} implementation that recovers an equivalent first-layer Value demixer and transfers it across requests without accessing the defender's secret parameters or fused dictionary.
\item We evaluate recovery, calibration cost, secret-epoch stability, closed-set value identification, and a simple refresh intervention, distinguishing raw-count experiments, legacy summaries, and unresolved properties.
\end{itemize}

\begin{figure*}[t]
\centering\includegraphics[width=\textwidth]{overview.pdf}
\caption{Attack mechanism. A legal repeat run makes first-layer Value rows identical. Fresh permutations move only the marker, producing a small orbit. The attacker recovers equivalent transforms from its own protected blocks, then matches demixed victim rows against a public-model dictionary. The diagram is schematic; no victim plaintext or server secret is an attack input.}
\label{fig:overview}
\end{figure*}
